% Gerado do mesmo CSV de componentes; não editar manualmente.

\noindent\textbf{\citeonline{Villegas2025_6916} -- Quantitativo não randomizado.}

\textbf{3.1 (CT):} A população de recrutamento e a estratégia de amostragem não estão estabelecidas nas seções examinadas. \textbf{3.2 (CT):} As medidas de resultado são descritas, mas o relato disponível não fundamenta suficientemente sua adequação. \textbf{3.3 (CT):} A completude dos dados de resultado e as perdas de participantes não estão estabelecidas. \textbf{3.4 (CT):} A alocação dos grupos e o controle inicial de fatores de confusão não estão suficientemente documentados. \textbf{3.5 (CT):} A introdução menciona um semestre acadêmico, os métodos relatam seis semanas e os resultados mostram progresso na oitava semana para estudantes de programação com alta interação; essa divergência não estabelece a adesão ou a cronologia do componente de matemática.

\noindent Localização: Páginas físicas 2 (450 estudantes e semestre letivo), 13–16 (delineamento, medidas e análise), 18–23 (resultados; semana 8 na p. 22) e 24–27 (testes e limites)\par

\noindent\textbf{\citeonline{Villegas2025_6916} -- Qualitativo.}

\textbf{1.1 (Y):} A abordagem qualitativa é pertinente para examinar as percepções de estudantes e educadores sobre o sistema de tutoria. \textbf{1.2 (Y):} Entrevistas estruturadas com estudantes e educadores, realizadas ao fim da experiência e com tópicos pertinentes explicitados, são adequadas para explorar as percepções sobre o sistema (página física 16). Detalhes de seleção e registro continuam como perguntas de esclarecimento, sem serem tratados automaticamente como inadequação da coleta. \textbf{1.3 (CT):} O relato menciona análise temática e temas, mas não explicita suficientemente a derivação analítica a partir das entrevistas (páginas físicas 16–17). \textbf{1.4 (CT):} As interpretações das entrevistas não são suficientemente vinculadas ao material qualitativo examinado (páginas físicas 16–17); médias e percentuais de satisfação nas pp. 24–25 são dados quantitativos e não resolvem esse vínculo. \textbf{1.5 (CT):} A coerência entre fontes qualitativas, coleta, análise e interpretação permanece pouco clara.

\noindent Localização: Páginas físicas 16–17 (Qualitative Analysis); 24–25 (satisfação quantitativa, não dados de entrevista)\par

\noindent\textbf{\citeonline{UniversityMathematics2026_6919} -- Quantitativo não randomizado.}

\textbf{3.1 (CT):} Os grupos pertencem a uma disciplina e estão separados por quatro anos. O artigo menciona amostragem por conglomerados, mas não define a população-alvo nem explica como as turmas de 2021 e 2025 foram escolhidas (páginas físicas 7–8). \textbf{3.2 (Y):} O questionário e os resultados da disciplina são descritos e analisados; a apreciação permanece provisória. \textbf{3.3 (CT):} A completude e o tratamento dos dados de resultado não estão suficientemente documentados nas seções examinadas. \textbf{3.4 (N):} Os grupos de controle e experimental de anos diferentes introduzem confundimento temporal reconhecido pelos autores; os testes de linha de base não controlam mudanças de equipe, contexto e uso de IA (páginas físicas 8 e 24). \textbf{3.5 (CT):} A implementação da intervenção e a adesão não estão suficientemente documentadas para este critério.

\noindent Localização: Páginas físicas 7–8 (§4.1.1), 12–14 (§4.3–4.4.1, tabela 4) e 24 (§6)\par

\noindent\textbf{\citeonline{UniversityMathematics2026_6919} -- Qualitativo.}

\textbf{1.1 (Y):} As entrevistas são adequadas para explorar as perspectivas dos estudantes sobre a reforma da disciplina. \textbf{1.2 (Y):} O artigo relata entrevistas semiestruturadas com 12 estudantes do grupo experimental, selecionados aleatoriamente logo após o curso (página física 18). \textbf{1.3 (CT):} As etapas analíticas utilizadas para derivar os temas não estão suficientemente detalhadas. \textbf{1.4 (Y):} A interpretação inclui relatos dos estudantes e é apoiada por citações apresentadas no artigo (página física 18). \textbf{1.5 (CT):} A relação entre amostragem qualitativa, coleta, análise e conclusões integradas requer verificação adicional.

\noindent Localização: Páginas físicas 18 (§4.4.8) e 24 (§6)\par

\noindent\textbf{\citeonline{Echeveria2025_6920} -- Quantitativo descritivo.}

\textbf{4.1 (Y):} Os registros de estudantes de uma escola são adequados à questão de classificação local descrita no estudo. \textbf{4.2 (CT):} A escola e seus estudantes delimitam uma população local; o MMAT não exige uma população-alvo mais abrangente. A correspondência dos 428 casos completos aos 633 recrutados não foi estabelecida, com anos letivos divergentes e sem comparação suficiente dos retidos e não retidos (páginas físicas 3 e 7). \textbf{4.3 (N):} Oito itens da escala usados para construir as categorias de mentalidade também permanecem entre os preditores selecionados, comprometendo a mensuração independente do alvo. \textbf{4.4 (CT):} O relato informa 633 participantes recrutados e 428 casos completos e cita os anos letivos 2023–2024 e 2022–2023 sem esclarecer a origem dos grupos nem se os casos ausentes produzem viés de não resposta. \textbf{4.5 (CT):} As descrições da divisão dos dados e da validação divergem e não permitem apreciar claramente a adequação da análise.

\noindent Localização: Páginas físicas 3/\allowbreak{}impressa 65 (§2.1–2.3), 6/\allowbreak{}68 (§2.6) e 7–8/\allowbreak{}69–70 (resultados)\par

\noindent\textbf{\citeonline{Echeveria2025_6920} -- Qualitativo.}

\textbf{1.1 (Y):} As entrevistas com professores são pertinentes aos objetivos declarados de contextualização e triangulação do estudo. \textbf{1.2 (Y):} O relato identifica entrevistas com cinco docentes de Matemática para explorar o uso do modelo e apresenta suas falas individuais (páginas físicas 3 e 11–12), constituindo coleta pertinente à questão de percepção docente. A seleção e o registro das respostas permanecem pouco detalhados. \textbf{1.3 (CT):} O procedimento analítico do material das entrevistas não está adequadamente descrito. \textbf{1.4 (Y):} As falas individuais dos docentes sustentam a interpretação descritiva de suas percepções sobre identificação, apoio e planejamento (páginas físicas 11–12). O julgamento não transforma expectativas dos entrevistados em efeito observado sobre aprendizagem. \textbf{1.5 (CT):} A coerência entre fontes, coleta, análise e interpretação qualitativas não pode ser concluída, pois o procedimento de análise das entrevistas permanece pouco explicitado (páginas físicas 3 e 11–13). A integração com dados quantitativos é examinada separadamente nos critérios de métodos mistos.

\noindent Localização: Páginas físicas 3/\allowbreak{}impressa 65 (§2.1 e coleta), 11–12/\allowbreak{}73–74 (falas docentes) e 13/\allowbreak{}75 (conclusões)\par

